\scanpage{52}
\mspage{50}
% archive sheet number: 51
principal \inserted{(\unsure{disons})} de représenter le foncteur $H$ \uline{sur les
catégories finies}. Et que les isom. obtenus
soient, bien entendu, \unsure{les bons|des bons} dans $\mathrm{Ind}(\mathrm{Cat}\,\mathrm{hot})$.
J'ai l'impression de ne pas avoir encore
axiomatisé à fond la notion même de
chemin, qui après tout a un \deleted{\unsure{signification}} \inserted{sens}
géométrique assez évident et important par
lui-même. Dire \unsure{peu|p.ex.} ce qu'il faut entendre
par ``\inserted{\uline{une}} théorie des chemins dans une catégorie''.
Une telle théorie \inserted{$\unsure{\mathcal T|\tau}$} associerait donc, à
$M,\Sigma,a,b$, un \deleted{\unsure{cat}} système inductif de
catégories $\mathrm{Ch}_{n\,\Sigma}^{\mathcal T}(M;a,b)$. Et les contraintes
géométriques naturelles à imposer à \deleted{\unsure{des}} des
telles théories, en vue \inserted{aussi} de leur utilisation
pour expliciter
\[
\mathrm{Hom}_{M(X)\Sigma_X^{-1}}(a_X,b_X)\simeq
\pi_0\underline{\mathrm{Hom}}\bigl(X,\underline{\underline{\mathrm{Ch}}}^{\mathcal T}(M;a,b)\bigr)
\]
donnerait un \inserted{syst. transitif d'}isom. \deleted{entre les} non seulmt
entre les ind-objets de $\mathrm{Cat}\,\mathrm{hot}$, mais \inserted{peut-être} aussi
entre les $\mathrm{Ch}_{\infty\,\Sigma}^{\mathcal T}(M;a,b)$, pour des
théories différentes. Si on dispose d'un
tel énoncé, la catégorie $\mathrm{Ch}_{\infty\,\Sigma}(M;a,b)$,
et plus \unsure{exactemt|exactement} ses images dans $\mathrm{Cat}\,\mathrm{hot}$ et $\mathrm{Hot}$,
